Table~\ref{tab:e2} gives the result, and the first thing it shows is how uneven learning is at this budget.
Skills are acquired in abrupt steps (Figure~\ref{fig:e1e2}b), so a run ends either near ceiling or near chance, and whether a given step arrives before training stops varies with the seed.
Even \skill{follow}, the simplest skill, was solved in only \ETwoPoolFollowTrainedSolved{} of the \ETwoPoolFollowTrainedRuns{} runs whose training set contained it.
Alongside each mean we therefore report how many seeds ended at 95\% or above, and we read the table by comparing those counts.

\paragraph{A skill left out of training was never acquired.}
Without \skill{induce} exercises, \skill{induce} accuracy is at or below its guessing level in every seed (\ETwoNoInduceInduceList\%); without \skill{compose} exercises, so is \skill{compose} (\ETwoNoComposeComposeTwoList\%).
Following and composing do not yield induction, and following and inducing do not yield composition.
Within TinyTextbook each skill is therefore necessary for itself, which is the minimality condition of Section~\ref{sec:ladder} at level G1.

\paragraph{Leaving out the simplest skill removed more than that skill.}
The models trained on \skill{compose} and \skill{induce}, but never on a single look-up, did not learn to look up (\ETwoNoFollowFollowList\% on \skill{follow}) and did not learn to compose either (\ETwoNoFollowComposeTwoList\%, which is the guessing level for two-step chains).
\skill{compose} was solved in \ETwoPoolComposeWithoutFollowSolved{} of \ETwoPoolComposeWithoutFollowRuns{} runs without \skill{follow} in the training set, against \ETwoPoolComposeWithFollowSolved{} of \ETwoPoolComposeWithFollowRuns{} runs with it (the full curriculum and the one without \skill{induce}; one-sided Fisher exact test, $p = \ETwoFisherP$).
A two-step exercise contains two look-ups, so in principle it supervises everything a one-step exercise does.
In practice the one-step exercise acted as a stepping stone without which the two-step skill was not found within the budget.
This also answers, for this benchmark, the worry raised in Section~\ref{sec:curriculum} that rule following cannot be removed separately from composition: it can, and removing it matters.

\paragraph{The full curriculum was not learned in every seed.}
With all three exercise types in training, \ETwoFullFollowSolved{} of \ETwoFullSeeds{} seeds solved \skill{follow}, \ETwoFullComposeTwoSolved{} solved \skill{compose} and \ETwoFullInduceSolved{} solved \skill{induce}; \ETwoFullAllSolved{} solved all three.
Across all conditions, the runs that fell short either made their first transition late, when the learning rate had already decayed too far for the remaining skills to follow, or never made it.
Nothing in our data explains which seeds these are.
The practical consequence is that a leave-one-out comparison at a fixed budget must be read through counts over seeds, as above, and Section~\ref{sec:protocol} lists the corrections we would make to the design.


\paragraph{Trained skills did not combine by themselves.}
\skill{mixed} exercises ask for a tabulated function and the induced function in one chain.
The models trained on the full curriculum score \ETwoFullMixedTwo\% on them (per seed: \ETwoFullMixedTwoList).
This is no better than guessing: when $g$ is the outer function, its input cannot be one of the two example keys, so the answer is one of three symbols that the examples determine, and when $g$ is the inner function the answer is uniform; a model that has induced $g$ and otherwise guesses therefore reaches $(1/3 + 1/5)/2 = 26.7\%$.
None of the \ETwoRunsAll{} runs in this experiment solved \skill{mixed}; the best reached \ETwoMixedMax\%.
Having each skill in the weights was not enough to obtain their combination, which contradicts hypothesis H5 at this scale.
\skill{compose} at depth 3 is likewise at its guessing level in every run, as expected from four layer passes trained on two-step chains; E3 examines depth directly.
